\subsection{NOTATIONS}

Let H be a hyperplane of E. Recall that E-H has two connected components, called the open half-spaces bounded by H. Their closures are called the closed half-spaces bounded by H. Let \(x, y \in E\) . Then x and y are said to be strictly on the same side of H if they are contained in the same open half-space bounded by H, or equivalently, if the closed segment with extremities x and y does not meet H; x and y are said to be on opposite sides of H if x belongs to one of the open half-spaces bounded by H and y to the other. If \(x \in E\) and \(t \in T\) , then x and t are said to be strictly on the same side of H if this is so for x and \(h + t\) for all \(h \in H\) .

Let A be a non-empty connected subset of E. For any hyperplane H of E that does not meet A, \(D_{H}(A)\) denotes the unique open half-space bounded by H that contains A. If N is a set of hyperplanes of E, none of which meet A, put

\[
D _ {\mathfrak {N}} (A) = \bigcap_ {H \in \mathfrak {N}} D _ {H} (A).\tag{1}
\]

If \(\mathrm{A}\) consists of a single point \(a\) , we write \(\mathrm{D_H}(a)\) and \(\mathrm{D}_{\mathfrak{N}}(a)\) instead of \(\mathrm{D_H}(\{a\})\) and \(\mathrm{D}_{\mathfrak{N}}(\{a\})\) .

